\begin{resumo}[Abstract]
\SingleSpacing\setlength{\parindent}{0pt}
\begin{otherlanguage*}{english}
This study characterizes mass-transport deposits and complexes on the northern slope of the Santos Basin to understand their relationships with topography associated with the Santos Channel and salt tectonics. Their spatial distribution, stratigraphic organization, morphometry, internal configuration, and confinement were investigated. The research used the BS-500 three-dimensional seismic volume, combining amplitude interpretation with the cosine of instantaneous phase attribute. Deposit top and basal surfaces were mapped, features associated with salt tectonics were characterized, and depressions interpreted as pockmarks were identified. Morphometric measurements were converted from time to depth using event-specific velocities estimated from well data and weighted by inverse distance, with quantification of the uncertainty components considered. Additionally, the horizontal-to-vertical spectral ratio technique was applied to ocean-bottom seismometer records. Ten events were distinguished, comprising three individual deposits and two complexes containing three and four events, respectively. Mapped areas ranged from 11.72 to 1,585.48 km², and estimated volumes ranged from 0.4440 to 131.4 km³. Topography associated with the Santos Channel and salt tectonics constrained deposit distribution, confinement, and thickening, while faults recorded post-depositional deformation in some events. Six internal seismic configurations were recognized, including widespread chaotic patterns and the coexistence of sectors with different degrees of organization. Stratigraphic relationships involving possible buried depressions and pockmarks indicated different stages of formation relative to deposit emplacement. The depth estimate derived from the horizontal-to-vertical spectral ratio showed local agreement with the contact between events D and C of complex 2, interpreted as a candidate interface responsible for the observed resonance response. The results demonstrate the influence of slope compartments on preserved deposit geometry and sediment accumulation, providing a comparative characterization of events at different scales.

\vspace{\onelineskip}
\noindent\textbf{Keywords}: mass-transport deposits; mass-complexes deposits; three-dimensional seismic reflection; morphometry; salt tectonics; Santos Basin.
\end{otherlanguage*}
\end{resumo}
